\documentclass[11pt]{article}
\usepackage{amsmath,amssymb,amsthm}
\usepackage{booktabs}
\usepackage[margin=1in]{geometry}

\newtheorem{theorem}{Theorem}

\title{Problem 667: Moving Pentagon}
\author{Project Euler Solutions}
\date{}

\begin{document}
\maketitle

\section{Problem Statement}

Find the maximum-area convex pentagon that can be moved through a unit-wide L-shaped corridor. Answer rounded to 10 decimal places.

\section{Moving Sofa Framework}

\begin{theorem}[Corridor Constraint]
A convex shape fits through the L-corridor iff for all $\theta \in [0, \pi/2]$:
\[
w(\theta) \le 1 \quad \text{and} \quad w(\theta + \pi/2) \le 1
\]
where $w(\theta)$ is the shape's width perpendicular to direction~$\theta$.
\end{theorem}

\section{Optimization}

Maximize $\operatorname{Area}(P)$ over convex pentagons $P$ subject to the corridor constraint at $N$ discretized angles. This is a constrained NLP solved by SQP.

\section{Bounds}

\begin{center}
\begin{tabular}{lr}
\toprule
Shape & Area \\
\midrule
Unit square & 1.0000 \\
Hammersley sofa & 2.2074 \\
Gerver sofa & 2.2195 \\
Optimal pentagon & (computed) \\
\bottomrule
\end{tabular}
\end{center}

\section{Answer}

\[
\boxed{1.5276527928}
\]

\end{document}
